% ==============================================================================
% 01_introduction.tex – Chapter 1
% ==============================================================================
\chapter{Introduction}
\label{ch:einleitung}

\section{Motivation}
\label{sec:motivation}
Eye tracking measures where a person is looking. This information is useful for
many applications: for controlling a computer without hands, for example for
people with physical disabilities, for usability studies of websites, for
learning platforms, and for research in psychology and medicine. Professional
eye trackers work with infrared light and special cameras. They are very
accurate, but expensive and not available everywhere.

An ordinary webcam, in contrast, is built into almost every laptop. If the gaze
could be estimated with a webcam, eye tracking would be possible for very many
people without extra hardware. But this is difficult. The images have a low
resolution, and the eyes are only a few pixels large. The light changes, the
head moves, many people wear glasses, and every person has a different eye and
face shape. A generic model trained on many persons therefore often fits a new
person badly.

\section{Problem Statement and Goal}
\label{sec:ziel}
One solution is \emph{personalization}: one takes a pretrained model and adapts
it to one person with a little data. For this, the person looks for a short time
at a point that moves across the screen (calibration). The calibration should be
short, however, because nobody wants to calibrate for several minutes before
every use. Personalization must therefore be \emph{data-efficient}.

The goal of this thesis is to study how well a pretrained vision model can be
adapted to one person with a short webcam calibration. The model is a ResNet-18
pretrained on ImageNet. Its backbone stays frozen, and only a small regression
head is trained. The head is first trained on the GazeCapture dataset and then
personalized with own recordings.

\section{Research Questions}
\label{sec:forschungsfragen}
The thesis answers four research questions (RQ):
\begin{description}[leftmargin=1.2cm, style=nextline]
  \item[RQ1] How much does a short personalization lower the error compared with
        a generic model and compared with a simple estimate that always gives
        the same point?
  \item[RQ2] Which training setting is suitable for the regression head
        (trainable layers, activation, loss function, learning rate, batch
        size)?
  \item[RQ3] How do the data split (random, temporal, two videos) and the design
        of the calibration (movement of the target, removal of transition
        frames) and a temporal smoothing of the predictions affect the
        result?
  \item[RQ4] What is the influence of the recording conditions (light,
        background, glasses, length, device or webcam)?
\end{description}

\section{Contributions}
\label{sec:beitrag}
The thesis makes the following contributions:
\begin{itemize}
  \item \textbf{Own recordings:} 27 calibration sessions of person~A (the
        author) with two laptop webcams, 18 of them evaluated, and two sessions
        of a second person (person~B) for comparison. The sessions differ in the movement of the target,
        light, background, glasses and length. The newest
        generation stores the phase of the target for every frame (\code{WAIT}
        or \code{MOVE}).
  \item \textbf{Reproducible pipeline:} fixed seeds, frozen backbone, stored
        features (feature cache) and stored predictions. One personalization
        run takes about two seconds on a CPU.
  \item \textbf{Corrected error measure:} a bug in the evaluation code was found
        and fixed. Baselines are always computed on the same data as the
        model.
  \item \textbf{Systematic study:} hyperparameters with five repeats and
        confidence intervals, a spatial analysis of the error, a comparison
        of data splits and calibration designs, and a study of the temporal
        smoothing of the predictions.
  \item \textbf{Web demo:} a demo in the browser shows personalization live:
        calibration, gaze tracking and evaluation (Chapter~6).
\end{itemize}

\section{Structure of the Thesis}
\label{sec:aufbau}
Chapter~\ref{ch:sota} reviews the state of the art. Chapter~\ref{ch:methodik}
explains the data, the model and the procedure. Chapter~\ref{ch:results} shows
the experiments and results. Chapter~\ref{ch:diskussion} answers the research
questions and names the limitations. Chapter~\ref{ch:fazit} summarizes the
thesis and gives an outlook.
